%%=============================================================================
%% Inleiding
%%=============================================================================

\chapter{\IfLanguageName{dutch}{Inleiding}{Introduction}}%
\label{ch:inleiding}
Webapplicaties vormen een essentieel onderdeel van moderne digitale dienstverlening. Naarmate organisaties steeds meer processen digitaliseren, groeit het belang van gebruiksvriendelijkheid als kwaliteitscriterium. Een webapplicatie die technisch correct functioneert maar moeilijk te gebruiken is, leidt tot frustratie, fouten en afhaken van gebruikers. Verschillende studies tonen aan dat ontwikkelteams nog steeds een sterke focus leggen op functionele en technische vereisten, terwijl usability‑aspecten vaak pas laat in het ontwikkelingsproces aandacht krijgen. Fernandez, Insfran en Abrahão stellen dat usability in requirements engineering structureel onderbelicht blijft, waardoor problemen pas zichtbaar worden wanneer de applicatie al grotendeels ontwikkeld is \autocite{Fernandez2011}. Dit wordt bevestigd door Valentim, Silva en Conte, die aantonen dat usability‑evaluaties in webprojecten doorgaans pas plaatsvinden tijdens of na de implementatiefase \autocite{Valentim2018}.
Het te laat identificeren van usability‑problemen heeft belangrijke gevolgen. Volgens Boehm stijgen de kosten van aanpassingen exponentieel naarmate men later in de ontwikkelingscyclus ingrijpt \autocite{Boehm1981}. Problemen die vroeg gedetecteerd worden, zijn relatief eenvoudig op te lossen, terwijl dezelfde problemen in een latere fase leiden tot herontwerp, extra ontwikkeltijd en hogere kosten. Bovendien heeft late detectie een negatieve impact op de gebruikerservaring, wat kan resulteren in lagere adoptie, verminderde efficiëntie en een daling van de algemene tevredenheid.
Deze problematiek is duidelijk zichtbaar in de casus die centraal staat in deze bachelorproef: een geanonimiseerde React‑gebaseerde webapplicatie voor online inschrijving. In deze applicatie haken gebruikers frequent af tijdens het doorlopen van een meerstapsformulier met client‑side validatie. Meerstapsformulieren zijn bijzonder gevoelig voor usability‑problemen, onder meer door cognitieve belasting, validatiefouten en onduidelijke navigatie. Norman benadrukt dat kleine ontwerpbeslissingen grote gevolgen kunnen hebben voor gebruikersfouten en frustratie \autocite{Norman2013}. Krug stelt dat gebruikers intuïtief moeten kunnen navigeren zonder onnodige cognitieve belasting \autocite{Krug2014}. Empirisch onderzoek toont aan dat multi‑step forms vaak leiden tot hogere drop‑off‑percentages wanneer de structuur, feedback of foutafhandeling niet optimaal is \autocite{BargasAvila2019}.
Telemetry‑tools zoals Microsoft Application Insights bieden nieuwe mogelijkheden om gebruikersgedrag objectief te meten. Deze tools registreren onder meer page views, click events, validatiefouten, funnels en time‑per‑step, waardoor een gedetailleerd beeld ontstaat van hoe gebruikers door een applicatie navigeren \autocite{MicrosoftAI,TelemetryUX2023}. Door deze data te combineren met traditionele usability‑metingen ontstaat een rijker en betrouwbaarder inzicht in de gebruiksvriendelijkheid van een applicatie.


\section{\IfLanguageName{dutch}{Probleemstelling}{Problem Statement}}%
\label{sec:probleemstelling}
Het centrale probleem in deze bachelorproef is dat gebruikers problemen ervaren bij het doorlopen van een meerstapsformulier in een React‑applicatie, terwijl er geen systematische manier bestaat om gebruikersgedrag vroegtijdig te meten. Hierdoor worden usability‑problemen pas laat zichtbaar, wanneer aanpassingen meer tijd en middelen vergen. De doelgroep van dit onderzoek bestaat uit IT‑professionals, in het bijzonder full‑stack developers en software engineers die verantwoordelijk zijn voor zowel de technische implementatie als de gebruikerservaring van webapplicaties. Voor deze doelgroep is het essentieel om te weten wanneer usability‑testing het meest effectief is: vroeg, laat of op beide momenten.

\section{\IfLanguageName{dutch}{Onderzoeksvraag}{Research question}}%
\label{sec:onderzoeksvraag}
De centrale onderzoeksvraag luidt:
Wat is het effect van usability‑tests die vroeg, laat of zowel vroeg als laat in het ontwikkelingsproces worden uitgevoerd op de gebruiksvriendelijkheid van een React‑applicatie, gemeten via Application Insights‑data, taakvoltooiingstijd, fouten en SUS‑scores?
Deze vraag wordt verder opgesplitst in deelvragen binnen het probleem‑ en oplossingsdomein.

\section{\IfLanguageName{dutch}{Onderzoeksdoelstelling}{Research objective}}%
\label{sec:onderzoeksdoelstelling}
Het doel van deze bachelorproef is om empirisch vast te stellen wanneer usability-testing de grootste impact heeft op de gebruiksvriendelijkheid van een webapplicatie. Het onderzoek levert:

\begin{itemize}
  \item wetenschappelijke meerwaarde door een lacune in de literatuur te vullen rond de timing van usability-tests en de combinatie met telemetry-data;
  \item praktische meerwaarde door ontwikkelteams concrete aanbevelingen te bieden over wanneer usability-testing het meest effectief is;
  \item technische meerwaarde door een proof-of-concept React-applicatie te ontwikkelen met geïntegreerde Application Insights.
\end{itemize}

\section{\IfLanguageName{dutch}{Opzet van deze bachelorproef}{Structure of this bachelor thesis}}%
\label{sec:opzet-bachelorproef}
In Hoofdstuk~\ref{ch:stand-van-zaken} wordt de bestaande literatuur besproken over usability, multi‑step forms, React‑applicaties en telemetry‑gebaseerde metingen.
In Hoofdstuk~\ref{ch:methodologie} wordt de methodologie toegelicht, inclusief de opzet van de usability‑tests, de dataverzameling en de statistische analyse.
In Hoofdstuk~\ref{ch:poc} wordt de proof-of-concept React-applicatie beschreven, inclusief de architectuur van de drie formulieren en de integratie met Microsoft Application Insights.
In Hoofdstuk~\ref{ch:vroege-usability-test} worden de resultaten van de vroege usability-test besproken voor Form~A en Form~B.
In Hoofdstuk~\ref{ch:ontwerpverbeteringen} worden de ontwerpverbeteringen beschreven die op basis van de vroege test werden doorgevoerd.
In Hoofdstuk~\ref{ch:late-usability-test} worden de resultaten van de late usability-test besproken voor Form~B en Form~C.
In Hoofdstuk~\ref{ch:vergelijkende-analyse} worden de resultaten van de drie testgroepen statistisch vergeleken via een eenweg-ANOVA en Tukey HSD post-hoc analyse.
In Hoofdstuk~\ref{ch:conclusie} worden de resultaten samengevat, wordt een antwoord geformuleerd op de onderzoeksvraag en worden aanbevelingen voor toekomstig onderzoek geformuleerd.
